\subsection{The Formula and Where It Comes From}

Let
\[
A=\begin{pmatrix}a&b\\c&d\end{pmatrix},
\qquad
A^{-1}=\begin{pmatrix}p&q\\r&s\end{pmatrix}.
\]
The condition $AA^{-1}=I_2$ produces two systems:
\[
\begin{cases}ap+br=1,\\cp+dr=0,
\end{cases}
\qquad
\begin{cases}aq+bs=0,\\cq+ds=1.
\end{cases}
\]
Thus finding an inverse already amounts to solving linear systems.

\begin{theorembox}[Inverse of a $2\times2$ matrix]
If
\[
\det A=ad-bc\neq0,
\]
then
\[
\boxed{A^{-1}=\frac1{ad-bc}
\begin{pmatrix}d&-b\\-c&a\end{pmatrix}}.
\]
If $ad-bc=0$, the inverse does not exist.
\end{theorembox}

The diagonal entries are exchanged, the off-diagonal signs are reversed, and the entire matrix is scaled by the reciprocal determinant.
